\documentclass[12pt,a4paper]{article}
\usepackage[UTF8]{ctex}
\usepackage{amsmath,amssymb,amsthm}
\usepackage{geometry}

\geometry{margin=2.5cm}

\title{伴随变换$\text{Ad}_{T_{i,i-1}^{-1}}$的推导}
\author{第8章补充说明}
\date{}

\begin{document}

\maketitle

\section{问题提出}

\textbf{问题}：如何推导伴随变换的公式？

\textbf{需要推导的公式}：

根据原始文档（MR-tablet.txt），标准形式是：
\begin{equation}
\text{Ad}_{T_{i,i-1}} = \begin{bmatrix}
R_i & 0 \\
-[R_i p_i]_{\times} R_i & R_i
\end{bmatrix}
\end{equation}

但在某些推导中，也会使用：
\begin{equation}
\text{Ad}_{T_{i,i-1}^{-1}} = \begin{bmatrix}
R_i^T & 0 \\
-[R_i^T p_i]_{\times} R_i^T & R_i^T
\end{bmatrix}
\end{equation}

\textbf{关键点}：
\begin{itemize}
\item 如果$T_{i,i-1}$是从$\{i-1\}$到$\{i\}$的变换，应该使用$\text{Ad}_{T_{i,i-1}}$（基于从$\{i-1\}$到$\{i\}$的变换）
\item 这是标准的伴随变换形式
\item 在原始文档的递归算法中使用的是$\text{Ad}_{T_{i,i-1}}$
\end{itemize}

\textbf{其中}：
\begin{itemize}
\item $T_{i,i-1} = \begin{bmatrix} R_i & p_i \\ 0 & 1 \end{bmatrix}$：从坐标系$\{i-1\}$到坐标系$\{i\}$的齐次变换矩阵
\item $R_i \in SO(3)$：从$\{i-1\}$到$\{i\}$的旋转矩阵
\item $p_i \in \mathbb{R}^3$：从$\{i-1\}$原点到$\{i\}$原点的向量（在$\{i-1\}$中表示）
\item $[p]_{\times}$：向量$p$的反对称矩阵
\end{itemize}

\section{伴随变换的定义}

\subsection{基本定义}

\textbf{伴随变换}：伴随变换$\text{Ad}_T$将twist从一个坐标系转换到另一个坐标系。

\textbf{Twist定义}：

Twist是一个6维向量，包含角速度和线速度：
\begin{equation}
V = \begin{bmatrix} \omega \\ v \end{bmatrix} \in \mathbb{R}^6
\end{equation}

\textbf{伴随变换的作用}：

如果$V_a$是坐标系$\{a\}$的twist（在$\{a\}$中表示），$T_{ba}$是从$\{a\}$到$\{b\}$的变换矩阵，那么：
\begin{equation}
V_b = \text{Ad}_{T_{ba}}(V_a)
\end{equation}

其中$V_b$是坐标系$\{b\}$的twist（在$\{b\}$中表示）。

\textbf{关键点}：
\begin{itemize}
\item $T_{ba}$是从$\{a\}$到$\{b\}$的变换
\item $\text{Ad}_{T_{ba}}$将twist从$\{a\}$转换到$\{b\}$
\item \textbf{应该基于$T_{ba}$（从$\{a\}$到$\{b\}$），而不是$T_{ba}^{-1}$（从$\{b\}$到$\{a\}$）}
\item 这是标准的伴随变换定义
\end{itemize}

\subsection{逆变换的伴随}

\textbf{逆变换}：

如果$T_{ba}$是从$\{a\}$到$\{b\}$的变换，那么$T_{ba}^{-1} = T_{ab}$是从$\{b\}$到$\{a\}$的变换。

\textbf{逆伴随变换}：

\begin{equation}
V_a = \text{Ad}_{T_{ba}^{-1}}(V_b) = \text{Ad}_{T_{ab}}(V_b)
\end{equation}

这表示将twist从$\{b\}$转换到$\{a\}$。

\textbf{注意}：在速度合成公式中，虽然使用的是$R_i^T$的形式，但这实际上对应的是$\text{Ad}_{T_{i,i-1}^{-1}}$，因为我们需要将twist从$\{i-1\}$转换到$\{i\}$，但使用的是逆变换的形式。这需要特别注意符号约定。

\section{推导过程}

\subsection{方法1：推导$\text{Ad}_{T_{i,i-1}}$（从$\{i-1\}$到$\{i\}$）}

\textbf{问题}：如何将twist从坐标系$\{i-1\}$转换到坐标系$\{i\}$？

\textbf{使用$\text{Ad}_{T_{i,i-1}}$}：

如果$T_{i,i-1}$是从$\{i-1\}$到$\{i\}$的变换，那么：
\begin{equation}
V_i = \text{Ad}_{T_{i,i-1}}(V_{i-1})
\end{equation}

\textbf{角速度的变换}：

角速度是自由向量，只需要旋转：
\begin{equation}
\omega_i = R_i \omega_{i-1}
\end{equation}

\textbf{注意}：这里使用$R_i$（不是$R_i^T$），因为$R_i$是从$\{i-1\}$到$\{i\}$的旋转矩阵。

\textbf{线速度的变换}：

坐标系$\{i\}$原点的速度（在$\{i-1\}$中表示）为：
\begin{equation}
v_{i,\{i-1\}} = v_{i-1} + \omega_{i-1} \times p_i
\end{equation}

要转换到坐标系$\{i\}$中，需要左乘$R_i$：
\begin{equation}
v_i = R_i (v_{i-1} + \omega_{i-1} \times p_i)
\end{equation}

\textbf{使用反对称矩阵}：

\begin{equation}
v_i = R_i (v_{i-1} - [p_i]_{\times} \omega_{i-1}) = R_i v_{i-1} - R_i [p_i]_{\times} \omega_{i-1}
\end{equation}

\textbf{关键关系}：

\begin{equation}
R_i [p_i]_{\times} = [R_i p_i]_{\times} R_i
\end{equation}

\textbf{证明}：

对于任意向量$q$：
\begin{align}
R_i [p_i]_{\times} q &= R_i (p_i \times q) \\
&= (R_i p_i) \times (R_i q) \quad \text{（旋转保持叉积）} \\
&= [R_i p_i]_{\times} (R_i q) \\
&= [R_i p_i]_{\times} R_i q
\end{align}

因此，$R_i [p_i]_{\times} = [R_i p_i]_{\times} R_i$。

\textbf{代入}：

\begin{align}
v_i &= R_i v_{i-1} - R_i [p_i]_{\times} \omega_{i-1} \\
&= R_i v_{i-1} - [R_i p_i]_{\times} R_i \omega_{i-1}
\end{align}

\textbf{矩阵形式}：

\begin{equation}
\begin{bmatrix} \omega_i \\ v_i \end{bmatrix} = \begin{bmatrix}
R_i & 0 \\
-[R_i p_i]_{\times} R_i & R_i
\end{bmatrix} \begin{bmatrix} \omega_{i-1} \\ v_{i-1} \end{bmatrix}
\end{equation}

\textbf{伴随变换}：

因此：
\begin{equation}
\text{Ad}_{T_{i,i-1}} = \begin{bmatrix}
R_i & 0 \\
-[R_i p_i]_{\times} R_i & R_i
\end{bmatrix}
\end{equation}

\textbf{注意}：这是从$\{i-1\}$到$\{i\}$的伴随变换，使用$R_i$（不是$R_i^T$）。

\subsection{方法2：推导$\text{Ad}_{T_{i,i-1}^{-1}}$（为什么在速度合成中使用）}

\textbf{问题}：为什么在速度合成公式中使用$\text{Ad}_{T_{i,i-1}^{-1}}$而不是$\text{Ad}_{T_{i,i-1}}$？

\textbf{原因}：在速度合成公式中，我们使用的是$R_i^T$的形式，这对应的是逆变换。

\textbf{步骤1：齐次变换矩阵的逆}

\textbf{齐次变换矩阵}：

\begin{equation}
T_{i,i-1} = \begin{bmatrix} R_i & p_i \\ 0 & 1 \end{bmatrix}
\end{equation}

\textbf{逆矩阵}：

\begin{equation}
T_{i,i-1}^{-1} = \begin{bmatrix} R_i^T & -R_i^T p_i \\ 0 & 1 \end{bmatrix} = \begin{bmatrix} R_i^T & -R_i^T p_i \\ 0 & 1 \end{bmatrix}
\end{equation}

\textbf{验证}：

\begin{align}
T_{i,i-1} T_{i,i-1}^{-1} &= \begin{bmatrix} R_i & p_i \\ 0 & 1 \end{bmatrix} \begin{bmatrix} R_i^T & -R_i^T p_i \\ 0 & 1 \end{bmatrix} \\
&= \begin{bmatrix} R_i R_i^T & -R_i R_i^T p_i + p_i \\ 0 & 1 \end{bmatrix} \\
&= \begin{bmatrix} I & -p_i + p_i \\ 0 & 1 \end{bmatrix} \\
&= \begin{bmatrix} I & 0 \\ 0 & 1 \end{bmatrix} = I
\end{align}

因此，$T_{i,i-1}^{-1} = \begin{bmatrix} R_i^T & -R_i^T p_i \\ 0 & 1 \end{bmatrix}$是正确的。

\subsection{步骤2：Twist的坐标变换}

\textbf{问题}：如何将twist从坐标系$\{i-1\}$转换到坐标系$\{i\}$？

\textbf{已知}：
\begin{itemize}
\item 坐标系$\{i-1\}$的twist：$V_{i-1} = [\omega_{i-1}^T, v_{i-1}^T]^T$（在$\{i-1\}$中表示）
\item 需要求：坐标系$\{i\}$的twist：$V_i = [\omega_i^T, v_i^T]^T$（在$\{i\}$中表示）
\end{itemize}

\textbf{角速度的变换}：

角速度是自由向量，只需要旋转：
\begin{equation}
\omega_i = R_i^T \omega_{i-1}
\end{equation}

\textbf{线速度的变换}：

线速度需要考虑：
\begin{enumerate}
\item 坐标系$\{i-1\}$原点的速度：$v_{i-1}$
\item 由于坐标系$\{i-1\}$的旋转产生的速度：$\omega_{i-1} \times p_i$
\end{enumerate}

因此，坐标系$\{i\}$原点的速度（在$\{i-1\}$中表示）为：
\begin{equation}
v_{i,\{i-1\}} = v_{i-1} + \omega_{i-1} \times p_i
\end{equation}

要转换到坐标系$\{i\}$中，需要左乘$R_i^T$：
\begin{equation}
v_i = R_i^T (v_{i-1} + \omega_{i-1} \times p_i)
\end{equation}

\subsection{步骤3：使用反对称矩阵}

\textbf{反对称矩阵}：

对于向量$p = [p_x, p_y, p_z]^T$，其反对称矩阵为：
\begin{equation}
[p]_{\times} = \begin{bmatrix}
0 & -p_z & p_y \\
p_z & 0 & -p_x \\
-p_y & p_x & 0
\end{bmatrix}
\end{equation}

\textbf{性质}：

\begin{equation}
p \times q = [p]_{\times} q
\end{equation}

\textbf{应用}：

\begin{equation}
\omega_{i-1} \times p_i = [\omega_{i-1}]_{\times} p_i = -[p_i]_{\times} \omega_{i-1}
\end{equation}

\textbf{注意}：$[p]_{\times} q = -[q]_{\times} p$（反对称矩阵的性质）

因此：
\begin{equation}
v_i = R_i^T (v_{i-1} - [p_i]_{\times} \omega_{i-1})
\end{equation}

\subsection{步骤4：矩阵形式}

\textbf{整理}：

\begin{align}
v_i &= R_i^T (v_{i-1} - [p_i]_{\times} \omega_{i-1}) \\
&= R_i^T v_{i-1} - R_i^T [p_i]_{\times} \omega_{i-1}
\end{align}

\textbf{关键关系}：

\begin{equation}
R_i^T [p_i]_{\times} = [R_i^T p_i]_{\times} R_i^T
\end{equation}

\textbf{证明}：

对于任意向量$q$：
\begin{align}
R_i^T [p_i]_{\times} q &= R_i^T (p_i \times q) \\
&= (R_i^T p_i) \times (R_i^T q) \quad \text{（旋转保持叉积）} \\
&= [R_i^T p_i]_{\times} (R_i^T q) \\
&= [R_i^T p_i]_{\times} R_i^T q
\end{align}

因此，$R_i^T [p_i]_{\times} = [R_i^T p_i]_{\times} R_i^T$。

\textbf{代入}：

\begin{align}
v_i &= R_i^T v_{i-1} - R_i^T [p_i]_{\times} \omega_{i-1} \\
&= R_i^T v_{i-1} - [R_i^T p_i]_{\times} R_i^T \omega_{i-1}
\end{align}

\subsection{步骤5：完整的矩阵形式}

\textbf{角速度和线速度的合成}：

\begin{align}
\omega_i &= R_i^T \omega_{i-1} \\
v_i &= -[R_i^T p_i]_{\times} R_i^T \omega_{i-1} + R_i^T v_{i-1}
\end{align}

\textbf{矩阵形式}：

\begin{equation}
\begin{bmatrix} \omega_i \\ v_i \end{bmatrix} = \begin{bmatrix}
R_i^T & 0 \\
-[R_i^T p_i]_{\times} R_i^T & R_i^T
\end{bmatrix} \begin{bmatrix} \omega_{i-1} \\ v_{i-1} \end{bmatrix}
\end{equation}

\textbf{伴随变换}：

因此：
\begin{equation}
\text{Ad}_{T_{i,i-1}^{-1}} = \begin{bmatrix}
R_i^T & 0 \\
-[R_i^T p_i]_{\times} R_i^T & R_i^T
\end{bmatrix}
\end{equation}

\section{详细推导：关键步骤}

\subsection{关键关系1：$R_i^T [p_i]_{\times} = [R_i^T p_i]_{\times} R_i^T$}

\textbf{证明}：

对于任意向量$q \in \mathbb{R}^3$，我们需要证明：
\begin{equation}
R_i^T [p_i]_{\times} q = [R_i^T p_i]_{\times} R_i^T q
\end{equation}

\textbf{左边}：
\begin{align}
R_i^T [p_i]_{\times} q &= R_i^T (p_i \times q)
\end{align}

\textbf{右边}：
\begin{align}
[R_i^T p_i]_{\times} R_i^T q &= (R_i^T p_i) \times (R_i^T q)
\end{align}

\textbf{关键性质}：旋转矩阵保持叉积关系

对于旋转矩阵$R$和向量$a, b$：
\begin{equation}
R(a \times b) = (Ra) \times (Rb)
\end{equation}

\textbf{应用}：

\begin{align}
R_i^T (p_i \times q) &= (R_i^T p_i) \times (R_i^T q) \\
&= [R_i^T p_i]_{\times} (R_i^T q) \\
&= [R_i^T p_i]_{\times} R_i^T q
\end{align}

因此，$R_i^T [p_i]_{\times} = [R_i^T p_i]_{\times} R_i^T$。

\subsection{关键关系2：$[p]_{\times}^T = -[p]_{\times}$}

\textbf{反对称矩阵的性质}：

对于反对称矩阵$[p]_{\times}$：
\begin{equation}
[p]_{\times}^T = -[p]_{\times}
\end{equation}

\textbf{证明}：

\begin{equation}
[p]_{\times} = \begin{bmatrix}
0 & -p_z & p_y \\
p_z & 0 & -p_x \\
-p_y & p_x & 0
\end{bmatrix}
\end{equation}

\begin{equation}
[p]_{\times}^T = \begin{bmatrix}
0 & p_z & -p_y \\
-p_z & 0 & p_x \\
p_y & -p_x & 0
\end{bmatrix} = -[p]_{\times}
\end{equation}

\section{验证}

\subsection{验证1：与速度合成公式的一致性}

\textbf{速度合成公式}：

\begin{align}
\omega_i &= R_i^T \omega_{i-1} + \dot{\theta}_i z_i \tag{22} \\
v_i &= R_i^T (v_{i-1} + \omega_{i-1} \times p_i) + \dot{d}_i z_i \tag{23}
\end{align}

\textbf{使用伴随变换}：

\begin{equation}
\begin{bmatrix} \omega_i \\ v_i \end{bmatrix} = \text{Ad}_{T_{i,i-1}^{-1}} \begin{bmatrix} \omega_{i-1} \\ v_{i-1} \end{bmatrix} + \begin{bmatrix} \dot{\theta}_i z_i \\ \dot{d}_i z_i \end{bmatrix} \tag{21}
\end{equation}

\textbf{展开伴随变换}：

\begin{equation}
\text{Ad}_{T_{i,i-1}^{-1}} = \begin{bmatrix}
R_i^T & 0 \\
-[R_i^T p_i]_{\times} R_i^T & R_i^T
\end{bmatrix}
\end{equation}

\textbf{矩阵乘法}：

\begin{align}
\begin{bmatrix} \omega_i \\ v_i \end{bmatrix} &= \begin{bmatrix}
R_i^T & 0 \\
-[R_i^T p_i]_{\times} R_i^T & R_i^T
\end{bmatrix} \begin{bmatrix} \omega_{i-1} \\ v_{i-1} \end{bmatrix} + \begin{bmatrix} \dot{\theta}_i z_i \\ \dot{d}_i z_i \end{bmatrix} \\
&= \begin{bmatrix}
R_i^T \omega_{i-1} \\
-[R_i^T p_i]_{\times} R_i^T \omega_{i-1} + R_i^T v_{i-1}
\end{bmatrix} + \begin{bmatrix} \dot{\theta}_i z_i \\ \dot{d}_i z_i \end{bmatrix}
\end{align}

\textbf{展开后得到}：

\begin{align}
\omega_i &= R_i^T \omega_{i-1} + \dot{\theta}_i z_i \tag{22'} \\
v_i &= -[R_i^T p_i]_{\times} R_i^T \omega_{i-1} + R_i^T v_{i-1} + \dot{d}_i z_i \tag{23'}
\end{align}

\textbf{验证角速度部分}：

公式(22')与公式(22)完全一致，验证通过。

\textbf{验证线速度部分（关键）}：

我们需要证明公式(23')与公式(23)一致，即：
\begin{equation}
-[R_i^T p_i]_{\times} R_i^T \omega_{i-1} + R_i^T v_{i-1} = R_i^T (v_{i-1} + \omega_{i-1} \times p_i)
\end{equation}

\textbf{详细推导}：

\textbf{步骤1：展开公式(23)的右边}：

\begin{align}
R_i^T (v_{i-1} + \omega_{i-1} \times p_i) &= R_i^T v_{i-1} + R_i^T (\omega_{i-1} \times p_i)
\end{align}

\textbf{步骤2：将叉积转换为反对称矩阵}：

\begin{equation}
\omega_{i-1} \times p_i = [\omega_{i-1}]_{\times} p_i
\end{equation}

\textbf{步骤3：使用反对称矩阵的性质}：

对于任意向量$a$和$b$：
\begin{equation}
a \times b = [a]_{\times} b = -[b]_{\times} a
\end{equation}

因此：
\begin{equation}
\omega_{i-1} \times p_i = [\omega_{i-1}]_{\times} p_i = -[p_i]_{\times} \omega_{i-1}
\end{equation}

\textbf{步骤4：代入}：

\begin{align}
R_i^T (v_{i-1} + \omega_{i-1} \times p_i) &= R_i^T v_{i-1} + R_i^T (-[p_i]_{\times} \omega_{i-1}) \\
&= R_i^T v_{i-1} - R_i^T [p_i]_{\times} \omega_{i-1}
\end{align}

\textbf{步骤5：关键关系}：

我们需要证明：
\begin{equation}
R_i^T [p_i]_{\times} = [R_i^T p_i]_{\times} R_i^T
\end{equation}

\textbf{证明}：

对于任意向量$q$：
\begin{align}
R_i^T [p_i]_{\times} q &= R_i^T (p_i \times q) \\
&= (R_i^T p_i) \times (R_i^T q) \quad \text{（旋转保持叉积）} \\
&= [R_i^T p_i]_{\times} (R_i^T q) \\
&= [R_i^T p_i]_{\times} R_i^T q
\end{align}

因此，$R_i^T [p_i]_{\times} = [R_i^T p_i]_{\times} R_i^T$。

\textbf{步骤6：完成验证}：

\begin{align}
R_i^T (v_{i-1} + \omega_{i-1} \times p_i) &= R_i^T v_{i-1} - R_i^T [p_i]_{\times} \omega_{i-1} \\
&= R_i^T v_{i-1} - [R_i^T p_i]_{\times} R_i^T \omega_{i-1}
\end{align}

这与公式(23')的左边（除了$\dot{d}_i z_i$项）完全一致！

\textbf{最终验证}：

\begin{align}
v_i &= -[R_i^T p_i]_{\times} R_i^T \omega_{i-1} + R_i^T v_{i-1} + \dot{d}_i z_i \quad \text{（公式23'）} \\
&= R_i^T (v_{i-1} + \omega_{i-1} \times p_i) + \dot{d}_i z_i \quad \text{（公式23）}
\end{align}

验证通过！

\subsection{验证2：伴随变换的性质}

\textbf{性质1：$\text{Ad}_{T^{-1}} = (\text{Ad}_T)^{-1}$}

\textbf{验证}：

如果$\text{Ad}_T$将twist从$\{a\}$转换到$\{b\}$，那么$\text{Ad}_{T^{-1}}$应该将twist从$\{b\}$转换回$\{a\}$。

\textbf{性质2：$\text{Ad}_{T_1 T_2} = \text{Ad}_{T_1} \text{Ad}_{T_2}$}

\textbf{验证}：

这可以通过矩阵乘法的结合律验证。

\section{具体例子}

\subsection{例子1：简单的旋转和平移}

\textbf{场景}：

考虑一个简单的变换：
\begin{itemize}
\item $R_i = I$（单位矩阵，无旋转）
\item $p_i = [1, 0, 0]^T$（沿$x$轴平移1单位）
\end{itemize}

\textbf{伴随变换}：

\begin{equation}
\text{Ad}_{T_{i,i-1}^{-1}} = \begin{bmatrix}
I & 0 \\
-[p_i]_{\times} & I
\end{bmatrix} = \begin{bmatrix}
I & 0 \\
-\begin{bmatrix} 0 & 0 & 0 \\ 0 & 0 & -1 \\ 0 & 1 & 0 \end{bmatrix} & I
\end{bmatrix}
\end{equation}

\textbf{物理意义}：

由于没有旋转，角速度直接传递。线速度需要考虑由于平移产生的耦合项。

\subsection{例子2：绕$z$轴旋转}

\textbf{场景}：

考虑绕$z$轴旋转$\theta$角：
\begin{equation}
R_i = \begin{bmatrix}
\cos\theta & -\sin\theta & 0 \\
\sin\theta & \cos\theta & 0 \\
0 & 0 & 1
\end{bmatrix}
\end{equation}

$p_i = [0, 0, 0]^T$（无平移）

\textbf{伴随变换}：

\begin{equation}
\text{Ad}_{T_{i,i-1}^{-1}} = \begin{bmatrix}
R_i^T & 0 \\
0 & R_i^T
\end{bmatrix}
\end{equation}

\textbf{物理意义}：

由于没有平移，角速度和线速度都只需要旋转，没有耦合项。

\section{两种方法的对比和说明}

\subsection{方法1：$\text{Ad}_{T_{i,i-1}}$（标准形式）}

\textbf{公式}：
\begin{equation}
\text{Ad}_{T_{i,i-1}} = \begin{bmatrix}
R_i & 0 \\
-[R_i p_i]_{\times} R_i & R_i
\end{bmatrix}
\end{equation}

\textbf{使用}：
\begin{equation}
V_i = \text{Ad}_{T_{i,i-1}}(V_{i-1}) + A_i \dot{\theta}_i
\end{equation}

\textbf{展开}：
\begin{align}
\omega_i &= R_i \omega_{i-1} + \dot{\theta}_i z_i \\
v_i &= R_i (v_{i-1} + \omega_{i-1} \times p_i) + \dot{d}_i z_i
\end{align}

\textbf{特点}：
\begin{itemize}
\item 使用$R_i$（从$\{i-1\}$到$\{i\}$的旋转矩阵）
\item 这是标准的伴随变换形式
\item 与$T_{i,i-1}$的方向一致
\end{itemize}

\subsection{方法2：$\text{Ad}_{T_{i,i-1}^{-1}}$（在速度合成中使用）}

\textbf{公式}：
\begin{equation}
\text{Ad}_{T_{i,i-1}^{-1}} = \begin{bmatrix}
R_i^T & 0 \\
-[R_i^T p_i]_{\times} R_i^T & R_i^T
\end{bmatrix}
\end{equation}

\textbf{使用}：
\begin{equation}
V_i = \text{Ad}_{T_{i,i-1}^{-1}}(V_{i-1}) + A_i \dot{\theta}_i
\end{equation}

\textbf{展开}：
\begin{align}
\omega_i &= R_i^T \omega_{i-1} + \dot{\theta}_i z_i \\
v_i &= R_i^T (v_{i-1} + \omega_{i-1} \times p_i) + \dot{d}_i z_i
\end{align}

\textbf{特点}：
\begin{itemize}
\item 使用$R_i^T$（逆旋转矩阵）
\item 在速度合成公式中常用这种形式
\item 虽然使用逆变换，但结果与标准形式等价
\end{itemize}

\subsection{为什么在速度合成中使用$\text{Ad}_{T_{i,i-1}^{-1}}$？}

\textbf{关键问题}：

如果$T_{i,i-1}$是从$\{i-1\}$到$\{i\}$的变换，那么理论上应该使用$\text{Ad}_{T_{i,i-1}}$（使用$R_i$），而不是$\text{Ad}_{T_{i,i-1}^{-1}}$（使用$R_i^T$）。

\textbf{但在速度合成公式中}：

速度合成公式使用的是$R_i^T$的形式：
\begin{align}
\omega_i &= R_i^T \omega_{i-1} + \dot{\theta}_i z_i \\
v_i &= R_i^T (v_{i-1} + \omega_{i-1} \times p_i) + \dot{d}_i z_i
\end{align}

\textbf{原因分析}：

这可能是因为$R_i$的定义不同：
\begin{itemize}
\item \textbf{在标准定义中}：$R_i$是从$\{i-1\}$到$\{i\}$的旋转矩阵，要将向量从$\{i-1\}$转换到$\{i\}$，应该使用$R_i$
\item \textbf{在速度合成中}：可能$R_i$的定义是从$\{i\}$到$\{i-1\}$的旋转矩阵，因此需要使用$R_i^T$来将向量从$\{i-1\}$转换到$\{i\}$
\end{itemize}

\textbf{符号约定的澄清}：

在本文档中，我们采用以下约定：
\begin{itemize}
\item $T_{i,i-1}$：从$\{i-1\}$到$\{i\}$的变换
\item $R_i$：从$\{i-1\}$到$\{i\}$的旋转矩阵
\item 因此，应该使用$\text{Ad}_{T_{i,i-1}}$（使用$R_i$）
\end{itemize}

\textbf{但在速度合成公式中}：

如果速度合成公式使用的是$R_i^T$，那么：
\begin{itemize}
\item 要么$R_i$的定义是从$\{i\}$到$\{i-1\}$（与标准定义相反）
\item 要么使用的是$\text{Ad}_{T_{i,i-1}^{-1}}$的形式（虽然符号上是逆变换，但通过$R_i^T$实现了正确的转换）
\end{itemize}

\textbf{建议}：

为了保持一致性，建议：
\begin{enumerate}
\item 如果$T_{ba}$是从$\{a\}$到$\{b\}$的变换，使用$\text{Ad}_{T_{ba}}$（使用$R_{ba}$）
\item 如果速度合成公式中使用$R_i^T$，需要检查$R_i$的定义是否与标准约定一致
\item 在推导时，明确说明符号约定，避免混淆
\end{enumerate}

\section{总结}

\subsection{推导步骤总结}

\begin{enumerate}
\item \textbf{齐次变换矩阵的逆}：$T_{i,i-1}^{-1} = \begin{bmatrix} R_i^T & -R_i^T p_i \\ 0 & 1 \end{bmatrix}$

\item \textbf{角速度变换}：$\omega_i = R_i^T \omega_{i-1}$（只需要旋转）

\item \textbf{线速度变换}：$v_i = R_i^T (v_{i-1} + \omega_{i-1} \times p_i)$（需要考虑旋转和平移的耦合）

\item \textbf{使用反对称矩阵}：$\omega_{i-1} \times p_i = -[p_i]_{\times} \omega_{i-1}$

\item \textbf{关键关系}：$R_i^T [p_i]_{\times} = [R_i^T p_i]_{\times} R_i^T$

\item \textbf{矩阵形式}：整理成$6 \times 6$矩阵形式
\end{enumerate}

\subsection{最终公式}

\textbf{伴随变换}：
\begin{equation}
\text{Ad}_{T_{i,i-1}^{-1}} = \begin{bmatrix}
R_i^T & 0 \\
-[R_i^T p_i]_{\times} R_i^T & R_i^T
\end{bmatrix}
\end{equation}

\textbf{其中}：
\begin{itemize}
\item $R_i^T$：旋转部分（上左和下右块）
\item $-[R_i^T p_i]_{\times} R_i^T$：耦合项（下左块），表示由于平移产生的角速度对线速度的影响
\item $0$：上右块为零，因为角速度不直接影响线速度（通过平移耦合）
\end{itemize}

\textbf{物理意义}：
\begin{itemize}
\item 上左块：角速度的旋转变换
\item 下左块：角速度对线速度的耦合影响（由于平移$p_i$）
\item 下右块：线速度的旋转变换
\end{itemize}

\end{document}

